% Part: methods
% Chapter: induction
% Section: structural-induction

\documentclass[../../../include/open-logic-section]{subfiles}

\begin{document}

\olfileid[pt]{mth}{ind}{sti}

\olsection{Indução Estrutural}

Até aqui usamos a indução para estabelecer resultados sobre todos os números
naturais. Mas um princípio correspondente pode ser usado diretamente para provar
resultados sobre todos os !!{element}s de um conjunto definido indutivamente.  Isso é
frequentemente chamado de indução \emph{estrutural}, porque depende da
estrutura dos objetos definidos indutivamente.

Em geral, uma definição indutiva é dada por (a)~uma lista de
!!{element}s ``iniciais'' do conjunto e (b)~uma lista de operações que
produzem novos !!{element}s do conjunto a partir de antigos. No caso dos termos
bonitos, por exemplo, os objetos iniciais são as letras. Temos apenas
uma operação: as operações são
\begin{align*}
  o(s_1, s_2) = & [s_1 \circ s_2]
\end{align*}
Você pode até pensar nos próprios números naturais~$\Nat$ como sendo
dados por uma definição indutiva: o objeto inicial é~$0$, e a
operação é a função sucessor~$x + 1$.

Para provar algo sobre todos os elementos de um conjunto definido
indutivamente, isto é, que todo !!{element} do conjunto tem uma
propriedade~$P$, devemos:
\begin{enumerate}
\item Provar que os objetos iniciais têm~$P$
\item Provar que, para cada operação~$o$, se os argumentos têm~$P$,
  o resultado também tem.
\end{enumerate}
Por exemplo, para provar algo sobre todos os termos bonitos, nós
provaríamos que é verdadeiro para todas as letras, e que é verdadeiro
para $[s_1 \circ s_2]$ desde que seja verdadeiro de $s_1$ e $s_2$
individualmente.

\begin{prop}
  O número de $[$ é igual ao número de $]$ em qualquer termo bonito~$t$.
\end{prop}

\begin{proof}
Usamos a indução estrutural.  Os termos bonitos são definidos indutivamente, com
as letras como objetos iniciais e a operação $o$ para construir novos
termos bonitos a partir de antigos.
\begin{enumerate}
\item A afirmação é verdadeira para toda letra, já que o número de $[$ em uma
  letra isolada é~$0$ e o número de $]$ nela também é~$0$.
\item Suponha que o número de $[$ em $s_1$ é igual ao número de $]$, e
  o mesmo é verdadeiro para $s_2$. O número de $[$ em $o(s_1, s_2)$, isto é, em
    $[s_1 \circ s_2]$, é a soma do número de $[$ em $s_1$ e
      $s_2$ mais um. O número de $]$ em $o(s_1, s_2)$ é a soma do número
    de $]$ em $s_1$ e $s_2$ mais um. Assim, o número de $[$ em $o(s_1, s_2)$
    é igual ao número de $]$ em $o(s_1,s_2)$.
\end{enumerate}
\end{proof}

\begin{prob}
  Prove, por indução estrutural, que nenhum termo bonito começa com~$]$.
\end{prob}

Vamos dar outra prova por indução estrutural: um segmento inicial próprio
de uma cadeia~$t$ de símbolos é qualquer cadeia~$s$ que concorda com
$t$ símbolo por símbolo, lida da esquerda, mas $t$ é mais longa. Assim,
por exemplo, $[a \circ {}$ é um segmento inicial próprio de $[a \circ b]$, mas
nem $[b \circ {}$ (discordam no segundo símbolo) nem $[a
\circ b]$ (têm o mesmo comprimento) o são.

\begin{prop}\ollabel{prop:initial}
  Todo segmento inicial próprio de um termo bonito~$t$ tem mais $[$ do que
    $]$.
\end{prop}

\begin{proof}
  Por indução em~$t$:
  \begin{enumerate}
  \item $t$ é uma letra isolada: Então $t$ não tem segmentos iniciais próprios.
  \item $t = [s_1 \circ s_2]$ para alguns termos bonitos $s_1$ e~$s_2$. Se $r$ é
    um segmento inicial próprio de $t$, há várias
    possibilidades:
    \begin{enumerate}
    \item $r$ é apenas $[$: Então $r$ tem um $[$ a mais do que~$]$.
    \item $r$ é $[r_1$ em que $r_1$ é um segmento inicial próprio
      de~$s_1$: Como $s_1$ é um termo bonito, por hipótese de indução, $r_1$
      tem mais $[$ do que $]$, e o mesmo é verdadeiro para $[r_1$.
    \item $r$ é $[s_1$ ou $[s_1 \circ {}$: Pelo resultado anterior, o
      número de $[$ e $]$ em~$s_1$ é igual; então o número de $[$
      em $[s_1$ ou $[s_1 \circ {}$ é um a mais do que o número de~$]$.
    \item $r$ é $[s_1 \circ r_2$ em que $r_2$ é um segmento inicial próprio
      de~$s_2$:  Por hipótese de indução, $r_2$ contém mais $[$ do que
      $]$. Pelo resultado anterior, o número de $[$ e de
      $]$ em~$s_1$ é igual. Então o número de $[$ em $[s_1 \circ r_2$
      é maior do que o número de~$]$.
    \item $r$ é $[s_1 \circ s_2$: Pelo resultado anterior, o número de
      $[$ e $]$ em $s_1$ é igual, e o mesmo para~$s_2$. Então há
      um $[$ a mais em $[s_1 \circ s_2$ do que há~$]$.
    \end{enumerate}
  \end{enumerate}
\end{proof}

\end{document}
